%%%%%%%%%%%%%%%%%%%%%%%%%%%%%%%%%%%%%%%%%%%%%%%%%%%%%%%%%%%%%%%%%%%%%
%% sec_intro.tex
%% Section 1: Introduction
%% RMG Periodic RT-TDDFT — Paper 2
%%%%%%%%%%%%%%%%%%%%%%%%%%%%%%%%%%%%%%%%%%%%%%%%%%%%%%%%%%%%%%%%%%%%%

\section{Introduction}
\label{sec:intro}

Real-time time-dependent density functional theory (RT-TDDFT) approach is a widely 
used for simulating electronic dynamics in molecules and materials 
under external perturbations.\cite{tddft-fundamentals,chem-rev-Li,Ullrich-2025-perspective}
%Real-time time-dependent density functional theory (RT-TDDFT) has
%established itself as one of the primary tools for simulating the
%electronic response of materials to external electromagnetic
%perturbations.\cite{tddft-fundamentals,chem-rev-Li}
Its applications span a broad range of phenomena, including optical
spectroscopy,\cite{lopata-2011,raman-govind-2013,Yabana-dichrosim-2011,Rubio-angular-2017,xray-herbert-2023}
charge transfer and photovoltaics,\cite{review-photovoltaics,chem-rev-Li,lvn-ct-2012}
nanoplasmonics,\cite{LSPR-review,wong-plasmonic-antennas-2017,tao-2021}
and spin and magnetization
dynamics.\cite{spin-rt-tddft-Li-2014,spin-rt-tddft-Li-2016,LvN-spin-Hod-Scuseria-2015}
In an earlier publication, we presented a scalable RT-TDDFT
implementation within the Real-space MultiGrid (\texttt{RMG}) code
for finite molecular systems, demonstrating simulations of up to
24,000 electrons on massively parallel
architectures.\cite{jakowski-rmg-tddft-2025}
The present work extends that implementation to periodic systems and
spin-polarized calculations, significantly broadening the range of
physically relevant problems accessible to \texttt{RMG}.

The transition from finite to periodic systems in RT-TDDFT requires
a corresponding transition in the treatment of the electromagnetic
perturbation.
In finite systems, the length gauge, in which a spatially uniform electric 
field couples through the position operator $\hat{\mathbf r}$, provides 
a natural description of the electric-dipole interaction. Under periodic 
boundary conditions, however, the associated scalar potential 
$\Phi_{\mathrm{ext}}(\mathbf r,t)=-\mathbf E(t)\cdot\mathbf r$ is not periodic;
enforcing periodicity produces a sawtooth potential with a discontinuity 
at the cell boundary.\cite{Luber-2023-position-operator}
Consequently, the usual position-based formulation cannot be applied directly 
to evaluate bulk polarization of an extended periodic systems, requiring 
a Berry-phase formulation.\cite{Resta-polarization-1994,KingSmith-Vanderbilt-polarization-1993}

Alternatively, the velocity gauge provides a convenient framework for periodic systems by representing 
the applied electric field through a spatially uniform, time-dependent vector potential
$\mathbf{A}(t)$, that preseres  the lattice periodicity of the Hamiltonian.
The length- and velocity-gauge formulations are related by a gauge transformation 
of the time-dependent Kohn--Sham equations,\cite{Mattiat-Luber-2022,furche-2001,ding-gauge-2011}
and requires careful treatment of the nonlocal part of
pseudopotentials.\cite{pickard-mauri-2001,starace-1971,starace-1973,Mattiat-Luber-2022}
The optical response is obtained from the induced macroscopic current density $\mathbf{J}(t)$
from which the optical conductivity
$\boldsymbol{\sigma}(\omega)$ and dielectric function
$\boldsymbol{\varepsilon}(\omega)$ are extracted by Fourier analysis.

RT-TDDFT for periodic systems has been implemented in several codes.
The \texttt{Octopus} code,\cite{octopus} based on real-space grids,
provides a comprehensive periodic RT-TDDFT implementation including
both length and velocity gauges and has served as the primary
validation reference for the present work.
Plane-wave codes including \texttt{Qbox}/\texttt{Qball},\cite{qball-2017,qball-II-2021}
\texttt{cp2k},\cite{cp2k-2020} and \texttt{Exciting}\cite{exciting-2026}
offer periodic RT-TDDFT with natural handling of Born-von K\'{a}rm\'{a}n
boundary conditions.
Atom-centered basis set codes such as \texttt{SIESTA}\cite{siesta-2007}
and \texttt{DFTB+}\cite{DFTB+} provide efficient implementations for
systems where localized orbitals are appropriate.
The \texttt{DFT-FE} code\cite{gavini-rt-tddft} offers an alternative
real-space approach based on finite elements.
%%\myblue{[Check: add any missing periodic RT-TDDFT codes before submission.
%%In particular: GPAW, Abinit, Yambo.]}

%The treatment of spin-polarized systems within RT-TDDFT introduces
%an additional degree of freedom: the electronic density is resolved
%into spin-up and spin-down components, and the time-dependent
%exchange-correlation potential depends on both. In the collinear
%approximation, the two spin channels evolve independently subject to
%a common self-consistent spin-dependent potential, and the optical
%response carries spin-specific signatures. This capability is
%essential for magnetic insulators such as \ch{CrI3}, where the
%spin-up and spin-down optical spectra differ markedly and encode
%information about the electronic and magnetic
%structure.\cite{spin-rt-tddft-Li-2014,LvN-spin-Hod-Scuseria-2015}
%The non-collinear spin extension, in which the electronic states are
%represented as two-component spinors, is beyond the scope of the
%present work and will be addressed in a future publication.

A distinctive feature of the \texttt{RMG} implementation is its
scalability. The real-space multigrid approach assigns different
regions of space to different processors, supporting massively
parallel calculations with near-linear scaling demonstrated in
earlier work on GPU-accelerated
architectures.\cite{jakowski-rmg-tddft-2025,rmg-dft-2024}
This scalability is particularly valuable for periodic
spin-polarized calculations, where sampling multiple k points
and propagating separate spin channels introduce additional
computational cost.
To establish the accuracy of the present implementation,
we validate its results against \texttt{Octopus} for
one-dimensional H$_2$ chains, two-dimensional hexagonal BN,
and three-dimensional Si, together with spin-polarized
calculations for \ch{CrI3}.
We then demonstrate its capacity to treat large periodic systems
through spin-polarized RT-TDDFT calculations of nitrogen-vacancy
(NV) centers in diamond, a prototypical point defect system of
substantial interest for quantum sensing and quantum information
applications.\cite{NV-Doherty-2013} 
Using supercells containing up to 4096 lattice sites
(4095 atoms and one vacancy), we investigate the dependence
of the defect optical response on supercell size and assess
the influence of interactions between periodic images.


The remainder of this paper is organized as follows.
Section~\ref{sec:methodology} presents the theoretical framework for
periodic spin-collinear RT-TDDFT, beginning with the interaction
of matter with electromagnetic fields,  the transformation
from length to velocity gauge, the current operator with nonlocal
pseudopotential contributions, density matrix propagation, and
the extraction of optical conductivity.
Section~\ref{sec:implementation} describes the implementation
in \texttt{RMG}, its principal approximations, and their physical
justification.
Section~\ref{sec:benchmarks} presents the validation benchmarks
and the large-supercell NV-center demonstration.
The paper concludes with a Summary and Outlook.


